%%%%%%%%%%%%%%%%%%%%%%%%%%%%%%%%%%%%%%%%%%%%%%%%%%%%%%
% COMPUTACIÓN CUÁNTICA: NOTACIÓN
%%%%%%%%%%%%%%%%%%%%%%%%%%%%%%%%%%%%%%%%%%%%%%%%%%%%%%

\usepackage{mathtools}
% babel en español vuelve "activos" caracteres como < y >, lo que rompe
% las flechas de TikZ (\draw[->]). Esta librería lo corrige.
\usetikzlibrary{babel, calc, decorations.pathreplacing}
% Además desactivamos los atajos < y > del español (solo sirven para escribir
% comillas «» con << >>), para que no interfieran con comparaciones en TikZ.
\AtBeginDocument{\shorthandoff{<>}}

% Color para las balotas azules de los diagramas de probabilidad
\definecolor{azulBalota}{RGB}{52,101,164}
% physics redefine \qty de siunitx: guardamos el de siunitx y lo restauramos
\NewCommandCopy\SIqty\qty
\usepackage{physics} % \ket, \bra, \braket, \ketbra, \expval, \abs, \norm, \tr, \comm
\RenewCommandCopy\qty\SIqty

% Conjuntos y espacios
\newcommand{\C}{\mathbb{C}}
\newcommand{\R}{\mathbb{R}}
\newcommand{\N}{\mathbb{N}}
\newcommand{\Hil}{\mathcal{H}}

%%%%%%%%%%%%%%%%%%%%%%%%%%%%%%%%%%%%%%%%%%%%%%%%%%%%%%
% ENCABEZADO DE CADA CLASE
%%%%%%%%%%%%%%%%%%%%%%%%%%%%%%%%%%%%%%%%%%%%%%%%%%%%%%
% Uso: \clase{fecha}{Título}. La fecha ya no se imprime: el documento
% está organizado por temas, como un libro.

\newcommand{\clase}[2]{\section{#2}}

%%%%%%%%%%%%%%%%%%%%%%%%%%%%%%%%%%%%%%%%%%%%%%%%%%%%%%
% CAJAS NUMERADAS (mismo estilo que tus cuadros)
%%%%%%%%%%%%%%%%%%%%%%%%%%%%%%%%%%%%%%%%%%%%%%%%%%%%%%

\tcbset{
  cajaUN/.style={
    enhanced, breakable,
    colback=white, colframe=#1, colbacktitle=#1, coltitle=white,
    fonttitle=\AncizarSansSemiBold,
    boxrule=0.8pt, arc=2mm, left=3mm, right=3mm, top=2mm, bottom=2mm
  }
}

% Uso: \begin{definicion}[Nombre opcional] ... \end{definicion}
\newtcolorbox[auto counter, number within=section]{definicion}[1][]{
  cajaUN=verdeOscuroUN,
  title={Definición~\thetcbcounter\ifx\relax#1\relax\else: #1\fi}
}

\newtcolorbox[use counter from=definicion]{teorema}[1][]{
  cajaUN=rojoUN,
  title={Teorema~\thetcbcounter\ifx\relax#1\relax\else: #1\fi}
}

\newtcolorbox[use counter from=definicion]{ejemplo}[1][]{
  cajaUN=grisOscuroUN,
  title={Ejemplo~\thetcbcounter\ifx\relax#1\relax\else: #1\fi}
}

% Material añadido desde la bibliografía (no dicho en clase)
% Uso: \begin{complemento}{Nielsen \& Chuang, §2.1} ... \end{complemento}
\newtcolorbox{complemento}[1]{
  enhanced, breakable,
  colback=grisClaroUN!12, colframe=grisOscuroUN,
  borderline west={3pt}{0pt}{grisOscuroUN},
  boxrule=0pt, arc=0mm, left=5mm, right=3mm, top=2mm, bottom=2mm,
  before upper={{\AncizarSansSemiBold\color{grisOscuroUN} Complemento}
                {\AncizarSansLightItalic\color{grisOscuroUN}\ (#1)}\par\smallskip}
}

%%%%%%%%%%%%%%%%%%%%%%%%%%%%%%%%%%%%%%%%%%%%%%%%%%%%%%
% CÓDIGO PYTHON / QISKIT
%%%%%%%%%%%%%%%%%%%%%%%%%%%%%%%%%%%%%%%%%%%%%%%%%%%%%%

\lstdefinestyle{Python}{
  style=Rcode,
  language=Python,
  morekeywords={QuantumCircuit, AerSimulator, Statevector, transpile, np}
}

%%%%%%%%%%%%%%%%%%%%%%%%%%%%%%%%%%%%%%%%%%%%%%%%%%%%%%
% OBSERVACIONES TÉCNICAS
% Comentario de motivación, alcance o conexión entre temas;
% no proviene de la transcripción literal de la clase.
% Uso: \begin{explicacion_tecnica}[Título opcional] ... \end{explicacion_tecnica}
%%%%%%%%%%%%%%%%%%%%%%%%%%%%%%%%%%%%%%%%%%%%%%%%%%%%%%

\newtcolorbox{explicacion_tecnica}[1][Observación]{
  enhanced, breakable,
  colback=grisClaroUN!8, colframe=verdeOscuroUN,
  borderline west={3pt}{0pt}{verdeOscuroUN},
  boxrule=0pt, arc=0mm, left=5mm, right=3mm, top=2mm, bottom=2mm,
  before upper={{\AncizarSansSemiBold\color{verdeOscuroUN} #1.}\ }
}
